% Tutorial, Chapter 6 -- More Commands and @Functions: Lessons 18 to 24
% (AN-17).  Figures: screens/tl18a..tl24.k123, each typed from an empty
% worksheet.  Sources (Law 1): os/i123/app/face_a.c (/Range Erase, /Range
% Label and its Left Right Center menu, /Range Value, /Range Trans,
% /Worksheet Column Set-Width, /Worksheet Erase and its No Yes menu),
% os/i123/core/range.c (Trans: the formulas of the FROM range arrive as
% values), os/i123/core/fmt.c (Date 1, 2 and 3; asterisks when the column
% is too narrow), os/i123/core/eval.c (@DATE, @AVG, @MAX, @MIN, @COUNT,
% @ROUND), spec/i123/era_shortfalls.md (Date 4, Date 5 and Time refused).
% Examples: examples/erase18, freeze20, dates22, funcs23, replayed.
% GATES: test-i123-range-menu test-i123-formula test-pubs-i123-examples test-pubs-screens test-pubs-i123
\Gates{test-i123-range-menu, test-i123-formula, test-pubs-i123-examples,
  test-pubs-screens, test-pubs-i123}
\chapter{More Commands and @Functions}

In this chapter you leave the budget and work with several small
worksheets, one for each new command.  Each lesson after the first begins
with an empty worksheet, so you can do the lessons in any order.  Start
Initech 123 again now; the worksheet is empty.

\section{Lesson 18: Erasing Cells}
\ix{/Range Erase}\ix{erasing}
Type a new entry over an old one and the old one is gone.  To remove an
entry and put nothing in its place, erase the cell.  First build a small
table of supplies.

\begin{Keystrokes}
\Do{Type}{\type{Pens}}{in A1, and press \kname{RIGHT}}
\Do{Type}{\type{12}}{in B1, press \kname{DOWN}, and then press
  \kname{LEFT}}
\Do{Type}{\type{Pads}}{in A2, press \kname{RIGHT}, type \type{7}, press
  \kname{DOWN} and \kname{LEFT}}
\Do{Type}{\type{Tape}}{in A3, press \kname{RIGHT}, type \type{4}, press
  \kname{DOWN} and \kname{LEFT}}
\Do{Type}{\type{Total}}{in A4, and press \kname{RIGHT}}
\Do{Type}{\type{@sum(b1..b3)}}{and press \kname{ENTER}}
\end{Keystrokes}
\FlairShot[fig:tl18a]{tl18a}{The table of supplies}
Pads are no longer stocked.  Erase the quantity in B2.

\begin{Keystrokes}
\Do{Press}{\kname{SLASH}}{and type \type{r}, \type{e} to select Range
  Erase}
\Do{Type}{\type{b2}}{the range to erase, and press \kname{ENTER}}
\end{Keystrokes}
The cell is empty and the total in B4 is recalculated without it
(Figure~\ref{fig:tl18b}).
\FlairShot[fig:tl18b]{tl18b}{B2 erased}
\LotExample{examples/erase18}
To erase several cells at once, give a range such as \texttt{A2..B2}.

\section{Lesson 19: Aligning Labels}
\ix{/Range Label}\ix{labels, aligning}
A label begins with a prefix character that tells Initech 123 how to place
it in its cell: an apostrophe for left, a quotation mark for right and a
caret for center.  A label you type without a prefix gets an apostrophe.
\slashcmd{/Range Label} changes the prefix of every label in a range at
once.

Erase the worksheet: press \kname{SLASH}, type \type{w}, \type{e},
\type{y}.  (Lesson~24 explains this command.)

\begin{Keystrokes}
\Do{Type}{\type{Pens}}{in A1, and press \kname{DOWN}}
\Do{Type}{\type{Pads}}{in A2, and press \kname{DOWN}}
\Do{Type}{\type{Tape}}{in A3, and press \kname{ENTER}}
\Do{Press}{\kname{SLASH}}{type \type{c}, type \type{a1..a3} and press
  \kname{ENTER}, type \type{b1} and press \kname{ENTER}}
\Do{Press}{\kname{SLASH}}{type \type{c}, type \type{a1..a3} and press
  \kname{ENTER}, type \type{c1} and press \kname{ENTER}}
\Do{Press}{\kname{SLASH}}{and type \type{r}, \type{l}, \type{r} to select
  Range Label Right}
\Do{Type}{\type{b1..b3}}{and press \kname{ENTER}}
\Do{Press}{\kname{SLASH}}{and type \type{r}, \type{l}, \type{c} to select
  Range Label Center}
\Do{Type}{\type{c1..c3}}{and press \kname{ENTER}}
\end{Keystrokes}
Column A is aligned at the left, column B at the right and column C in the
center (Figure~\ref{fig:tl19}).  The third choice on the menu,
\slashcmd{/Range Label Left}, puts the apostrophe back.
\FlairShot[fig:tl19]{tl19}{The same labels aligned left, right and center}

\section{Lesson 20: Freezing Values}
\ix{/Range Value}\ix{values, freezing}
A formula changes whenever the cells it refers to change.  Sometimes you
want to keep the answer it gives today.  \slashcmd{/Range Value} copies a
range and turns every formula in the copy into the number it displays.

Erase the worksheet, and then build a small sum.

\begin{Keystrokes}
\Do{Type}{\type{12.5}}{in A1, and press \kname{DOWN}}
\Do{Type}{\type{7.5}}{in A2, and press \kname{DOWN}}
\Do{Type}{\type{@sum(a1..a2)}}{in A3, and press \kname{ENTER}}
\Do{Press}{\kname{SLASH}}{and type \type{r}, \type{v} to select Range
  Value}
\Do{Type}{\type{a3}}{the range to copy FROM, and press \kname{ENTER}}
\Do{Type}{\type{c3}}{the cell to copy TO, and press \kname{ENTER}}
\Do{Press}{\fkey{Goto}{F5}}{and type \type{a1}, then press \kname{ENTER}}
\Do{Type}{\type{20}}{and press \kname{ENTER}}
\Do{Press}{\fkey{Goto}{F5}}{and type \type{c3}, then press \kname{ENTER}}
\end{Keystrokes}
The total in A3 follows the change in A1, but C3 keeps the value it had
when you copied it; the control panel shows a number, not a formula
(Figure~\ref{fig:tl20}).
\FlairShot[fig:tl20]{tl20}{A3 follows A1; C3 is a frozen value}
\LotExample{examples/freeze20}

\section{Lesson 21: Exchanging Rows and Columns}
\ix{/Range Trans}\ix{transposing}
\slashcmd{/Range Trans} copies a range so that its rows become columns and
its columns become rows.  Erase the worksheet, and type a table of three
quarters.

\begin{Keystrokes}
\Do{Type}{\type{"Q1}}{in A1, press \kname{RIGHT}; \type{"Q2},
  \kname{RIGHT}; \type{"Q3}, \kname{ENTER}}
\Do{Press}{\fkey{Goto}{F5}}{and type \type{a2}, then press \kname{ENTER}}
\Do{Type}{\type{10}}{press \kname{RIGHT}; \type{20}, \kname{RIGHT};
  \type{30}, \kname{ENTER}}
\Do{Press}{\kname{SLASH}}{and type \type{r}, \type{t} to select Range
  Trans}
\Do{Type}{\type{a1..c2}}{the range to copy FROM, and press \kname{ENTER}}
\Do{Type}{\type{a4}}{the upper left cell of the copy, and press
  \kname{ENTER}}
\end{Keystrokes}
The two rows of three cells become two columns of three cells
(Figure~\ref{fig:tl21}).  The original stays where it is.
\FlairShot[fig:tl21]{tl21}{The table copied with rows and columns exchanged}
\begin{Note}
Like \slashcmd{/Range Value}, \slashcmd{/Range Trans} copies the displayed
result of a formula, not the formula itself.
\end{Note}

\section{Lesson 22: Entering Dates}
\ix{dates}\ix{"@DATE}\ix{/Range Format Date}
Initech 123 keeps a date as a \term{date serial number}: a count of days that
begins in 1900.  The function
\texttt{@DATE(\Arg{year},\Arg{month},\Arg{day})} gives the serial number
of a date, with the last two digits of the year.  Because dates are
numbers, you can subtract one from another.  Erase the worksheet and
enter two dates and the difference between them.

\begin{Keystrokes}
\Do{Type}{\type{@date(92,6,15)}}{in A1, and press \kname{DOWN}}
\Do{Type}{\type{@date(92,12,31)}}{in A2, and press \kname{DOWN}}
\Do{Type}{\type{+a2-a1}}{in A3, and press \kname{ENTER}}
\end{Keystrokes}
A1 and A2 show serial numbers, and A3 shows the number of days between the
dates (Figure~\ref{fig:tl22a}).
\FlairShot[fig:tl22a]{tl22a}{Dates as serial numbers}
To display a serial number as a date, format the cells.  The format
DD-MMM-YY needs ten characters, so first widen column A.

\begin{Keystrokes}
\Do{Press}{\kname{SLASH}}{and type \type{w}, \type{c}, \type{s} to select
  Worksheet Column Set-Width}
\Do{Type}{\type{10}}{and press \kname{ENTER}}
\Do{Press}{\kname{SLASH}}{and type \type{r}, \type{f}, \type{d} to select
  Range Format Date}
\Do{Type}{\type{1}}{to select the format DD-MMM-YY}
\Do{Type}{\type{a1..a2}}{and press \kname{ENTER}}
\end{Keystrokes}
\FlairShot[fig:tl22b]{tl22b}{The same dates in the format DD-MMM-YY}
In a column that is too narrow the date appears as a row of asterisks;
widen the column and the date returns.  Formats 2 and 3 show only the day
and month (DD-MMM) or only the month and year (MMM-YY).  The cell keeps its
serial number whichever format you choose.
\LotExample{examples/dates22}

\section{Lesson 23: More @Functions}
\ix{"@AVG}\ix{"@MAX}\ix{"@MIN}\ix{"@COUNT}\ix{"@ROUND}
You used \fn{@SUM} in Lesson~5 and \fn{@IF} in Lesson~15.  Five more
functions complete a short list of the ones you will use most:

\begin{LotTable}{p{30mm}p{98mm}}
\LotHead{\textbf{Function} & \textbf{Result}}
\fn{@COUNT} & How many cells in the range are not blank\\
\fn{@AVG} & The average of the range\\
\fn{@MAX} & The largest value in the range\\
\fn{@MIN} & The smallest value in the range\\
\fn{@ROUND} & A value rounded to the number of decimal places you give;
  \texttt{@ROUND(\Arg{value},0)} rounds to a whole number\LotEnd
\end{LotTable}
Erase the worksheet, type five numbers in A1..A5, and ask the five
functions about them in C1..C5.

\begin{Keystrokes}
\Do{Type}{\type{8}}{in A1, press \kname{DOWN}; then \type{11}, \type{9} and
  \type{14}, each followed by \kname{DOWN}}
\Do{Type}{\type{6.5}}{in A5, and press \kname{ENTER}}
\Do{Press}{\fkey{Goto}{F5}}{and type \type{c1}, then press \kname{ENTER}}
\Do{Type}{\type{@count(a1..a5)}}{and press \kname{DOWN}}
\Do{Type}{\type{@avg(a1..a5)}}{and press \kname{DOWN}}
\Do{Type}{\type{@max(a1..a5)}}{and press \kname{DOWN}}
\Do{Type}{\type{@min(a1..a5)}}{and press \kname{DOWN}}
\Do{Type}{\type{@round(c2,0)}}{and press \kname{ENTER}}
\end{Keystrokes}
Column C now holds the count, the average, the largest value, the smallest
value and the average rounded to a whole number
(Figure~\ref{fig:tl23}).
\FlairShot[fig:tl23]{tl23}{Five functions asked about the same column}
\LotExample{examples/funcs23}
\begin{Caution}
\fn{@COUNT} counts a label as well as a number, and \fn{@AVG} treats a
label in its range as a zero.  Keep headings out of the ranges you give
them.
\end{Caution}

\section{Lesson 24: Starting Over}
\ix{/Worksheet Erase}
\slashcmd{/Worksheet Erase} clears the whole worksheet from memory, so you
can start a new one without ending Initech 123.  It does not change any
file on disk.  Because it cannot be undone, Initech 123 asks first.

\begin{Keystrokes}
\Do{Press}{\kname{SLASH}}{and type \type{w}, \type{e} to select Worksheet
  Erase}
\end{Keystrokes}
The menu offers \texttt{No} and \texttt{Yes}
(Figure~\ref{fig:tl24}).  \texttt{No} leaves the worksheet alone and returns
to \ind{READY} mode.
\FlairShot[fig:tl24]{tl24}{Worksheet Erase asks before it clears}
\begin{Keystrokes}
\Do{Type}{\type{y}}{to select Yes}
\end{Keystrokes}
The worksheet is empty and the pointer is in A1.  The unsaved entries are
gone; save a worksheet you want to keep with \slashcmd{/File Save} first.

This completes the \doctitle{Tutorial}.  The \doctitle{Initech 123
Reference} describes every command you have used, and all the others.
